\subsection{Aplicación y comparación de clasificadores}

Selecciona conjuntos de datos para entrenar distintos modelos de
clasificación variando la configuración de sus hiperparámetros.

\textbf{Conjuntos de datos propuestos:}

\begin{itemize}
    \item \textbf{Clasificación Multimodelo para Cáncer de Mama}
          (\href{https://archive.ics.uci.edu/dataset/17/breast+cancer+wisconsin+diagnostic}{\textit{Breast Cancer Wisconsin Diagnostic}}):
          conjunto de datos biomédicos que contiene características numéricas
          derivadas de imágenes digitales de aspiraciones con aguja fina (FNA)
          de masas mamarias, utilizadas para clasificar tumores como benignos
          o malignos.
    \item \textbf{Detección Temprana de Parkinson, voz y biomarcadores acústicos}
          (\href{https://archive.ics.uci.edu/dataset/189/parkinsons+telemonitoring}{\textit{Parkinsons Telemonitoring}}):
          conjunto diseñado para apoyar la investigación en el diagnóstico de
          la enfermedad de Parkinson. Contiene medidas biomédicas de voz
          obtenidas de personas sanas y pacientes con la enfermedad, con el
          objetivo de diferenciar ambos grupos mediante técnicas de aprendizaje
          automático.
    \item \textbf{Clasificación de Cardiopatía}
          (\href{https://archive.ics.uci.edu/dataset/45/heart+disease}{\textit{Heart Disease}}):
          recopila información de pacientes para determinar la probabilidad de
          enfermedad cardíaca, codificada en una variable objetivo que indica
          la presencia o ausencia de la condición. Entre las variables
          predictoras se incluyen edad, sexo, presión arterial en reposo, nivel
          de colesterol, frecuencia cardíaca máxima y resultados de pruebas
          electrocardiográficas.
    \item \textbf{Verificador de síntomas de COVID-19}
          (\href{https://www.kaggle.com/datasets/iamhungundji/covid19-symptoms-checker}{\textit{COVID Symptom Classification Dataset}}):
          los datos identifican si una persona padece o no la enfermedad del
          coronavirus, basándose en algunos síntomas estándar predefinidos.
          Dichos síntomas se basan en las directrices de la OMS y del
          Ministerio de Salud y Bienestar Familiar de la India.
\end{itemize}

\begin{enumerate}
    \item \textbf{Carga y análisis de datos:}
          \begin{enumerate}
              \item Cargar el conjunto de datos.
              \item Visualización de la distribución de clases.
              \item Análisis de valores faltantes (porcentaje por columna).
          \end{enumerate}

    \item \textbf{Preprocesamiento y limpieza:}
          \begin{enumerate}
              \item Separación de tipos de variables (columnas numéricas y categóricas).
              \item Verificar y aplicar limpieza de datos (faltantes,
                    \textit{outliers}, caracteres inválidos, etc.).
              \item Estandarización de características (escaladores).
              \item Separación en 70\% entrenamiento y 30\% prueba.
              \item Manejo del desbalance de clases sobre el conjunto de
                    entrenamiento (por ejemplo SMOTE, ADASYN o sobremuestreo
                    aleatorio).
              \item Visualización de la distribución.
          \end{enumerate}

    \item \textbf{Reducción de dimensionalidad por PCA:}
          \begin{enumerate}
              \item Aplicar PCA a los datos balanceados.
              \item Cálculo de la varianza acumulada.
              \item Encontrar el número de componentes para el 95\% de varianza.
              \item Gráfico del número de componentes principales y varianza
                    acumulada.
          \end{enumerate}

    \item \textbf{Definición de modelos, esquema de trabajo y búsqueda en
              cuadrícula:}
          \begin{enumerate}
              \item Desarrollar un esquema de funciones (\textit{pipelines})
                    que integren el preprocesamiento y el clasificador.
              \item Búsqueda en cuadrícula (\textit{grid search}) para encontrar
                    los mejores hiperparámetros.
          \end{enumerate}

    \item \textbf{Configuración de modelos y cuadro de hiperparámetros:}
          \begin{enumerate}
              \item \textbf{KNN}:
                    \begin{itemize}
                        \item vecinos (1, 3, 5, 7, 9).
                        \item pesos (\textit{weights: ['uniform', 'distance']}).
                        \item distancia (\textit{metric: ['euclidean', 'manhattan']}).
                    \end{itemize}
              \item \textbf{Árbol de decisión}:
                    \begin{itemize}
                        \item profundidad (\textit{[3, 5, 7, None]}).
                        \item muestras mínimas por división (\textit{[2, 5, 10]}).
                        \item criterio (\textit{['gini', 'entropy']}).
                    \end{itemize}
              \item \textbf{SVM}:
                    \begin{itemize}
                        \item kernel (\textit{['linear', 'rbf']}).
                        \item Regularización $C$ (0.1, 1, 10).
                        \item Gamma (\textit{['scale', 'auto', 0.01, 0.1]}).
                    \end{itemize}
          \end{enumerate}

    \item \textbf{Entrenamiento, evaluación y comparación de resultados:}
          \begin{enumerate}
              \item Mejores métricas de rendimiento sobre conjunto de prueba
                    por conjunto de datos.
              \item Matriz de confusión de los modelos entrenados con las
                    mejores métricas.
              \item Fronteras de decisión superponiendo datos de entrenamiento
                    y prueba.
              \item Árbol entrenado con los mejores parámetros.
              \item Comparación de exactitud entre modelos.
              \item ¿Qué hiperparámetro fue más relevante por modelo para cada
                    conjunto de datos?
          \end{enumerate}
\end{enumerate}
